% Generated draft: requires mathematical and source review.
\documentclass[11pt]{article}
\usepackage[margin=25mm]{geometry}
\usepackage{fontspec}
\setmainfont{Latin Modern Roman}
\usepackage{amsmath,amssymb,mathtools,unicode-math}
\setmathfont{Latin Modern Math}
\usepackage{enumitem,needspace,xurl,hyperref}
\setlength{\parindent}{0pt}
\setlength{\parskip}{5pt}
\setlength{\emergencystretch}{3em}
\newcommand{\N}{\mathbb{N}}
\newcommand{\Z}{\mathbb{Z}}
\newcommand{\Q}{\mathbb{Q}}
\newcommand{\R}{\mathbb{R}}
\newcommand{\C}{\mathbb{C}}
\newcommand{\F}{\mathbb{F}}
\newcommand{\A}{\mathbb{A}}
\newcommand{\PP}{\mathbb{P}}
\newcommand{\E}{\mathbb{E}}
\newcommand{\eps}{\varepsilon}
\newcommand{\dd}{\,\mathrm{d}}
\begin{document}
\section*{UnsolvedMath ID: 701009}
\textbf{Source problem code:} catalogue entry
\subsection{Definitions and mathematical statement}
Let \(F_n\) be a free group of rank \(n\), and let \(G_n=\operatorname{Aut}(F_n)\) be its group of automorphisms. The source equips \(G_n\) with a norm \(\lVert\cdot\rVert\) defined in Kourovka Notebook Problem 14.23. That definition is not included in the supplied record; the statement below is therefore conditional on that specified norm and does not substitute another norm or word length. Two elements \(\alpha,\beta\in G_n\) are conjugate if there exists \(\eta\in G_n\) with \(\eta^{-1}\alpha\eta=\beta\). A recursive function here means a total computable function.

Does there exist a recursive function \(f:\mathbb{N}\times\mathbb{N}\to\mathbb{N}\) such that
\[
\forall\alpha,\beta\in G_n,\qquad
\bigl(\exists\eta\in G_n:\eta^{-1}\alpha\eta=\beta\bigr)
\Longrightarrow
\bigl(\exists\gamma\in G_n:\gamma^{-1}\alpha\gamma=\beta
\ \text{and}\ 
\lVert\gamma\rVert\leq f(\lVert\alpha\rVert,\lVert\beta\rVert)\bigr)?
\]
The source does not explicitly specify the admissible ranks \(n\), or whether \(n\) is fixed before choosing \(f\), allowing dependence on \(n\), or whether one function must work for all ranks. These qualifications remain unresolved.
\subsection{Short English statement}
For conjugate automorphisms of a free group, is there a computable bound on the norm of some conjugating automorphism that depends only on the norms of the two given automorphisms? The intended norm is defined in Problem 14.23, which was not supplied, and the required dependence on the rank is unclear.
\subsection{Sources}
\begin{itemize}
\item Local UnsolvedMath registry: permanent ID \(701009\), exact source code \texttt{catalogue entry}, title \emph{Kourovka Notebook Problem 14.24}; supplied upstream revision \texttt{d371ee917ac9647fbb06cd6fc529deec6f3dda82}.
\item Supplied attribution: O. V. Bogopolski, proposed in 1999. Source: \emph{The Kourovka Notebook (Unsolved Problems in Group Theory)}, 14th issue (1999), Problem 14.24; the record also identifies the 21st issue edition (2026), \texttt{arXiv:1401.0300v48}. The supplied collection description names E. I. Khukhro and V. D. Mazurov as editors.
\item The statement refers to Problem 14.23 for its norm; the text of that problem was not supplied. The registry labels Problem 14.24 as open. This label and the bibliographic information are supplied source data, not an independently verified literature assessment.
\end{itemize}
\end{document}
